\documentclass[11pt]{article}       % set main text size
\usepackage[letterpaper,                % set paper size to letterpaper. change to a4paper for resumes outside of North America
top=0.5in,                          % specify top page margin
bottom=0.5in,                       % specify bottom page margin
left=0.5in,                         % specify left page margin
right=0.5in]{geometry}              % specify right page margin
                       
\usepackage{XCharter}               % set font. comment this line out if you want to use the default LaTeX font Computer Modern
\usepackage[T1]{fontenc}            % output encoding
\usepackage[utf8]{inputenc}         % input encoding
\usepackage{enumitem}               % enable lists for bullet points: itemize and \item
\usepackage[hidelinks]{hyperref}    % format hyperlinks
\usepackage{titlesec}               % enable section title customization
\usepackage{comment}
\raggedright                        % disable text justification
\pagestyle{empty}                   % disable page numbering

% ensure PDF output will be all-Unicode and machine-readable
\input{glyphtounicode}
\pdfgentounicode=1

% format section headings: bolding, size, white space above and below
\titleformat{\section}{\bfseries\large}{}{0pt}{}[\vspace{1pt}\titlerule\vspace{-6.5pt}]

% format bullet points: size, white space above and below, white space between bullets
\renewcommand\labelitemi{$\vcenter{\hbox{\small$\bullet$}}$}
\setlist[itemize]{itemsep=-2pt, leftmargin=12pt, topsep=7pt} %%% Test various topsep values to fix vertical spacing errors

% resume starts here
\begin{document}

% name
\centerline{\Huge Bradley Pascoe}

\vspace{5pt}

% contact information
\centerline{\href{mailto:bradley.pascoe@unisq.edu.au}{bradley.pascoe@unisq.edu.au} | \href{https://www.linkedin.com/in/bradleypascoe}{www.linkedin.com/in/bradleypascoe}}

\vspace{-10pt}

% skills section

% \vspace{-6.5pt}

% experience section
\section*{Experience}

\subsection*{Research}
\vspace{-6pt}
\textbf{Research Fellow} {University of Southern Queensland} -- Springfield, QLD \hfill Jun 2025 -- Present \\
\vspace{-9pt}
\begin{itemize}
    \item Redesigned research code for vehicle codesign optimisation into modular, tested Python libraries, enabling reuse across multiple studies and improving reproducibility, numerical methods, and performance.
    \item Developed a Python library for numerical computation of vehicle radiation spectra, accelerating core algorithms with Cython for high-performance evaluation.
\end{itemize}

\textbf{Doctoral Candidate, Research Associate} {University of Sydney} -- Sydney, NSW \hfill Mar 2021 -- Mar 2025 \\
\vspace{-9pt}
\begin{itemize}
    \item Investigated the effects of strain on turbulent mixing layers by implementing and validating advanced numerical methods into a Fortran codebase for the solution of partial differential equations. 
    \item Utilised high-performance computing infrastructure using code developed for OpenMP and MPI parallelisations to simulate and analyse terabytes of data, aided by Bash and Python scripts to organise jobs and files.
    % \item Automated visualisations of simulated data with Python APIs, greatly reducing the time required and reducing the memory usage by 20-75\%.
    \item Automated visualisations of simulated data with Python APIs, reducing time and memory usage by 20-75\%.
    \item Derived and calibrated reduced order models, achieving accurate predictions of fluid flow quantities for over 10,000$\times$ less computational cost compared to simulating the three-dimensional fluid flow.
\end{itemize}

\textbf{Research Engineer} {Defence Science and Technology Group} -- Melbourne, VIC \hfill Feb 2020 -- Jan 2021 \\
\vspace{-9pt}
\begin{itemize}
  \item Enhanced the pre- and post-processing capabilities of multi-dimensional fatigue spectrum and strain data by developing and optimising a library in R, including introducing version control and unit tests.
  \item Designed programs to read and evaluate data from fatigue testing systems, allowing for continuous verification of experimental regulatory compliance.
\end{itemize}
\vspace{-16pt}
\subsection*{Teaching}
\vspace{-6pt}
\subsubsection*{Academic Tutor/Demonstrator}
\vspace{-6pt}
Engineering Mechanics (AMME1802), University of Sydney \hfill 2024\\
Engineering Methods (AMME3060), University of Sydney \hfill 2024\\
Applied Fluid Dynamics and Turbulence (AMME5292), University of Sydney \hfill 2024\\
Linear Systems and Control (MMAN3200), University of NSW \hfill 2018, 2019, 2021, 2022\\
Engineering Mechanics (ENGG1300), University of NSW \hfill 2018, 2019 \\
Engineering Mechanics 2 (MMAN2300), University of NSW \hfill 2018, 2019 \\
Mechanics of Solids (MMAN2400), University of NSW \hfill 2018, 2019 \\
Physics 1A (PHYS1121), University of NSW \hfill 2018, 2019 \\

\vspace{-6pt}
\subsubsection*{Mentoring}
\vspace{-6pt}
Supervised an engineering undergraduate thesis project (AMME4112) \hfill 2023
\begin{comment}


% projects section
\section*{Projects}
\textbf{Project Title} \hfill \href{https://github.com/matiassingers/awesome-readme}{github.com/name/project} \\
\vspace{-9pt}
\begin{itemize}
  \item The more work experience you have, the less relevant outside-work projects tend to become
  \item If you have something that really stands out, consider listing it here
\end{itemize}

\textbf{Project Title} \hfill \href{https://mitcommlab.mit.edu/meche/commkit/portfolio/}{project.com} \\
\vspace{-9pt}
\begin{itemize}
  \item Only list real projects, not mandatory school projects or trivial tutorial projects found online 
  \item Something that someone uses to solve a problem
  \item Something that has users (can be just you, as long as you use it often) and is actively maintained and not just rotting in a GitHub repo, never to see a PR for the rest of its life
\end{itemize}

\textbf{Project Title} \hfill \href{https://www.hardwareishard.net/portfolio-database}{blog.com/projectname} \\
\vspace{-9pt}
\begin{itemize}
  \item Resume checklist: \url{https://old.reddit.com/r/EngineeringResumes/wiki/checklist}
  \item Google Docs version of this template: \url{https://docs.google.com/document/d/1MBvhATv8y-ESORopRoLSZ3f3HjkM_Qa_f8fIHAEqgnI/edit}
\end{itemize}

% alternate formatting for projects section
% \section*{Projects}
% \begin{itemize}
%   \item \textbf{QuantSoftware Toolkit:} Open source Python library for data analysis and machine learning for finance
%   \item \textbf{GitHub Visualization:} Data visualization of Git Log data using D3 to analyze project trends over time
%   \item \textbf{Recommendation System:} Music and movie recommender systems using collaborative filtering on public datasets
%   \item \textbf{Mac Setup:} Book that gives step by step instructions on setting up developer environment on Mac OS
% \end{itemize}
\end{comment}

% \vspace{-18.5pt}
%\vspace{-18.5pt}
% \vspace{-6.5pt}

% education section
\section*{Education}
\textbf{University of Sydney} -- PhD in Engineering \hfill Dec 2024 \\
\textbf{University of New South Wales} -- B.Eng.(Hons) in Aerospace Engineering \hfill Jan 2021 \\
\textbf{University of New South Wales} -- B.Sc. in Physics \hfill Jan 2021

\vspace{-6.5pt}

\section*{Skills}
% \textbf{CAD:} Siemens NX, CATIA V5, SolidWorks \\
% \textbf{Analysis:} Thermal Desktop, Abaqus, LS-DYNA, STAR-CCM+\\
\textbf{Languages:} Bash, C++, Fortran, MATLAB, Python, R, SQL\\
\textbf{Tools:} Git, LaTeX, MPI, OpenMP \\
\textbf{Software:} ANSYS Fluent, CATIA, Paraview, SolidWorks, Tecplot

\vspace{-6.5pt}
\section*{Publications}
\begin{itemize}
    \item \textbf{B. Pascoe}, M. Groom, D.L. Youngs, and B. Thornber (2026), Impact of transverse strain on linear, transitional and self-similar turbulent mixing layers, Journal of Fluid Mechanics, 1030, A61
    \item \textbf{B. Pascoe}, M.Groom, B. Thornber (2026), Late-time growth of an inhomogenous, turbulent mixing layer subjected to transient compression, Physica D: Nonlinear Phenomena, 488, 135056
    \item \textbf{B. Pascoe}, M.Groom, B. Thornber (2025), Turbulence modelling of mixing layer under anisotropic strain, Physical Review Fluids, 10, 064609
    \item \textbf{B. Pascoe}, M. Groom, D.L. Youngs, and B. Thornber (2024), Impact of axial strain on linear, transitional and self-similar turbulent mixing layers, Journal of Fluid Mechanics, 999, A5
\end{itemize}

\vspace{-16.5pt}
\section*{Presentations}
\begin{itemize}
\item \textit{Shock-driven compressible turbulent mixing in planar and convergence geometries}, University of Oxford (2024)
\item \textit{The impact of strain-rates on linear, transitional and self-similar turbulent mixing layers}, First Light Fusion (2024)
\item \textit{Reynolds-averaged Navier-Stokes modelling for the Richtmyer-Meshkov instability under anisotropic strain}, 18th International Worksop on the Physics of Compressible Turbulent Mixing (2024)
\item \textit{Modelling the effects of convergence on the Richtmyer-Meshkov instability in planar geometry}, American Physical Society Division of Fluid Dynamics Annual Meeting (2023)
\item \textit{Modelling the Richtmyer-Meshkov instability under strain}, International High Performance Computing Summer School (2023)
\item \textit{The Richtmyer-Meshkov instability under axial strain}, 23rd Australasian Fluid Mechanics Conference (2022)
\item \textit{The effect of axial strain on the Richtmyer-Meshkov instability}, 17th International Workshop on the Physics of Compressible Turbulent Mixing (2022)
\end{itemize}

\vspace{-16.5pt}
\section*{Grants}
\begin{itemize}
    \item Very High-Order Solver Frameworks for Compressible Turbulent Mixing, UK Research and Innovation GPU eCSE, £127,992 (2025)
\end{itemize}

\vspace{-16.5pt}
\section*{Honours and Awards}
\begin{itemize}
    \item University of Sydney HDR Publication Award - Mechanical/Thermofluids (2024)
    \item International High Performance Computing Summer School attendee (2023)
    \item Science Vacation Research Scholarship (2018)
    \item Future Physicist International Summer Camp attendee (2018)
\end{itemize}

\vspace{-16.5pt}
\section{Service}

\textbf{Journal/Conference Reviewer}\\
Journal of Fluid Mechanics, Physica D: Nonlinear Phenomena, 24th Australasian Fluid Mechanics Conference\\

\vspace{9pt}
\textbf{Elected Student Representative Council Member}\\
Australasian Fluid Mechanics Society \hfill 2023 -- 2024
\begin{itemize}
    \item Chair of the HDR/ECR subcommittee and the Emerging Leaders subcommittee.
\end{itemize}

% \vspace{6pt}
\textbf{Conference Administration Lead}\\
University of Sydney \hfill Sep 2022 -- Jan 2023
\begin{itemize}
  \item Prepared the delegate materials, childcare services, and catering, and managed a team of 30 volunteers for a four-day conference with 400+ delegates, helping the conference to achieve a rating of 4.57/5 for organisation.
  \item Designed and distributed the post-conference survey and compiled the results into a report, obtaining a record number of responses for the post-conference survey.
\end{itemize}

\end{document}
